\textbf{COLUNI 2019 -- Questão 09}\par\medskip Leia o diálogo a seguir retirado do filme \textit{O carteiro e o poeta}.

\par

Dom Pablo: Aqui nesta ilha, o mar. Tanto mar. Ele transborda de tempos a tempos. Diz que sim, diz que não, então não. No azul, na espuma, num galope diz que não, então não. Não pode estar tranquilo. “Chamo-te mar”, repete batendo na pedra sem convencê-la. Então, com sete línguas verdes de sete tigres verdes, de sete cães verdes, de sete mares verdes ele acaricia, beija, molha e bate no peito repetindo seu nome.\par
Dom Pablo: Então. Que te parece?\par
Mário: É Estranho.\par
Dom Pablo: Como assim, estranho? És um crítico severo.\par
Mário: Não, não o seu poema. Estranho… estranho… É como me senti enquanto estava a recitar.\par
Dom Pablo: E como se sentiu?\par
Mário: Não sei. As palavras iam para frente e para trás.\par
Dom Pablo: Como o mar?\par
Mário: Exatamente, como o mar.\par
Dom Pablo: Isso é ritmo.\par
Mário: Na verdade, senti-me mareado.\par
Dom Pablo: Mareado.\par
Mário: Não sei explicar. Senti-me como um barco balançando em volta de tuas palavras.\par
Dom Pablo: Como um barco balançando em volta das minhas palavras?\par
Dom Pablo: Sabes o que acabastes de fazer, Mário?\par
Mário: o quê?\par
Dom Pablo: Uma metáfora.\par
Mário: Não!\par
Dom Pablo: Sim, sim!\par
Mário: Mesmo? Mas não vale porque não tive a intenção.\par
Dom Pablo: As imagens nascem espontaneamente.

\par

Assinale a alternativa em que o verso NÃO apresenta o mesmo recurso linguístico que aquele utilizado por Mário nesse diálogo do filme:\begin{enumerate}
\item “O nosso mestre foi saindo da sala, meio rindo a chorar.” (BARROS, Manoel de. \textbf{Memórias inventadas}. Rio de Janeiro:
      Alfaguara, 2018, p. 44.)
\item “Até as pedras da rua choravam.” (BARROS, Manoel de. \textbf{Memórias inventadas}. Rio de Janeiro: Alfaguara, 2018, p. 64)
\item “[...] Esse menino vai passar a vida enfiando água no espeto!” (BARROS, Manoel de. \textbf{Memórias inventadas}. Rio de Janeiro:
      Alfaguara, 2018, p. 62)
\item “Pintei sem lápis a Manhã de pernas abertas para o Sol.” (BARROS, Manoel de. \textbf{Memórias inventadas}. Rio de Janeiro:
      Alfaguara, 2018, p.37)\end{enumerate}
